\section{Producing a contracting block}\label{sec:contraction}

The reduction will remove a small group of indices from the synthesis problem.
For its geometric series to converge rapidly, the corresponding diagonal block
of the unitary must have norm bounded away from one.  We now obtain such a block
using only permutations, diagonal signs, and Walsh--Hadamard transforms.  All
these transformations have short programmable-slot implementations.

We first record the matrix estimate needed to choose the signs.  A
\emph{Rademacher variable} takes the values $1$ and $-1$ with equal probability.
The following elementary moment argument is a finite-dimensional form of the
noncommutative Khintchine estimate; see Tropp~\cite[Sections~6--7]{Tropp2018}
for the Gaussian integration-by-parts and trace-product method.  We include
the proof with the constants used below.

\begin{lemma}[A matrix sign estimate]\label{contract:moment}
Let $L_1,\ldots,L_q$ be Hermitian $h\times h$ matrices, and let
$\varepsilon_1,\ldots,\varepsilon_q$ be independent Rademacher variables.
Set $v=\norm{\sum_i L_i^2}$.  For every integer $a\geq1$,
\begin{equation}\label{contract:sign-bound}
 \E\norm{\sum_{i=1}^q\varepsilon_iL_i}
 \leq 2h^{1/(2a)}\sqrt{(2a-1)v}.
\end{equation}
\end{lemma}

\begin{proof}
We first replace the signs by independent standard real Gaussians $g_i$ and
write $X=\sum_i g_iL_i$.  Gaussian integration by parts, applied to this matrix
polynomial, gives
\begin{equation}\label{contract:ibp}
 \E\Tr X^{2a}
 =\sum_{i=1}^q\sum_{l=0}^{2a-2}
   \E\Tr\bigl(L_iX^lL_iX^{2a-2-l}\bigr).
\end{equation}
Indeed, write the first factor of $X^{2a}$ as $\sum_i g_iL_i$, use
$\E[g_iP(g)]=\E[\partial_iP(g)]$, and differentiate each of the remaining
$2a-1$ factors in turn.

Put $m=2a-2$.  For a fixed Hermitian $X$, diagonalize it with real eigenvalues
$\lambda_j$.  In this basis, for $m>0$,
\begin{align*}
 \Tr(L_iX^lL_iX^{m-l})
 &=\sum_{j,k}|(L_i)_{jk}|^2\lambda_k^l\lambda_j^{m-l}\\
 &\leq\sum_{j,k}|(L_i)_{jk}|^2
   \left(\frac lm|\lambda_k|^m+
              \frac{m-l}{m}|\lambda_j|^m\right)
 =\Tr(L_i^2X^m).
\end{align*}
The inequality is the weighted arithmetic--geometric mean inequality after
taking absolute values.  The last equality uses the symmetry
$|(L_i)_{jk}|^2=|(L_i)_{kj}|^2$ and the evenness of $m$.
For $m=0$ the same assertion is equality.  Since $X^m$ is positive semidefinite
and $\sum_iL_i^2\preceq vI$, equation~\eqref{contract:ibp} implies
\[
 \E\Tr X^{2a}\leq(2a-1)v\,\E\Tr X^{2a-2}.
\]
Iterating from $\Tr I=h$ yields
\[
 \E\Tr X^{2a}
 \leq h(2a-1)!!\,v^a
 \leq h(2a-1)^a v^a.
\]
Here $(2a-1)!!=1\cdot3\cdots(2a-1)$.  Therefore
\begin{equation}\label{contract:gaussian-bound}
 \E\norm X
 \leq\bigl(\E\norm X^{2a}\bigr)^{1/(2a)}
 \leq h^{1/(2a)}\sqrt{(2a-1)v}.
\end{equation}

To return to signs, write $g_i=\varepsilon_i t_i$, where the $t_i=|g_i|$ are
independent of the signs and have common mean $\mu=\sqrt{2/\pi}$.  Conditional
Jensen gives
\[
 \mu\,\norm{\sum_i\varepsilon_iL_i}
 \leq\E_t\norm{\sum_i\varepsilon_i t_iL_i}.
\]
Average over the signs, apply~\eqref{contract:gaussian-bound}, and use
$\mu^{-1}<2$.
\end{proof}

Fix throughout the construction
\begin{equation}\label{contract:parameters}
 M=4096(n+1)^2,\qquad r=\frac34.
\end{equation}
For an integer $2\leq d\leq2^n$, let $q$ be the largest power of two not
exceeding $d$, and define
\begin{equation}\label{contract:orders}
 s=\max\{1,\lfloor q/M\rfloor\},\qquad k=d-s.
\end{equation}
We call the first $k$ indices \emph{ports} and the last $s$ indices
\emph{internal indices}.  Both groups are nonempty.

\begin{proposition}[A block of uniformly small norm]\label{contract:block}
For every unitary $U$ on $\C^d$, there are unitaries $P_{\mathrm L}$ and
$P_{\mathrm R}$ on $\C^d$ such that
\begin{equation}\label{contract:block-form}
 V=P_{\mathrm L}UP_{\mathrm R}
 =\begin{pmatrix}A&B\\C&D\end{pmatrix},
 \qquad \norm D\leq r,
\end{equation}
where the block sizes are $k,s$ from~\eqref{contract:orders}.
The transformations $P_{\mathrm L},P_{\mathrm R}$ and their inverses have
implementations using $O(n)$ programmable slots and $O(n)$ clean working
qubits, with a structure determined by $n,d$.  The same structure works
simultaneously for any family $U_\lambda$, using the label $\lambda$ as an
additional control.
\end{proposition}

\begin{proof}
If $s=1$, the last column of $U$ has an entry of modulus at most $1/\sqrt d$.
A row permutation moves this entry to the last position.  The resulting
one-dimensional block satisfies $\norm D\leq1/\sqrt d<3/4$, and the permutation
has the claimed slot implementation by Lemma~\ref{slot:permutation}.

Suppose $s>1$.  Let $M_0$ be the last-by-last $q\times q$ submatrix of $U$;
it is a contraction.  Set
$H_q=H^{\otimes\log_2q}$, and let $W$ consist of its last $s$ columns.  We
will choose diagonal sign matrices $E_1,E_2$ on $\C^q$ and use
$H_qE_1$ on the left and $E_2H_q$ on the right, acting as the identity on the
remaining $d-q$ indices.  The resulting internal block is
\begin{equation}\label{contract:D-signs}
 D=W^*E_1C',\qquad C'=M_0E_2W.
\end{equation}
We choose $E_2$ first to bound the entries of $C'$, then choose $E_1$ to
bound the operator norm of $D$.

For independent signs in $E_2$, each entry of $C'$ is a sign sum
$\sum_j\varepsilon_j\alpha_j$ with $\sum_j|\alpha_j|^2\leq1/q$.
For completeness, a real sign sum with squared coefficient sum at most $1/q$
satisfies
\[
 \E\exp\left(t\sum_j\varepsilon_j a_j\right)
 =\prod_j\cosh(ta_j)\leq\exp(t^2/(2q)).
\]
Markov's inequality, optimized over $t$, bounds its two-sided tail at $u$ by
$2e^{-qu^2/2}$.  Applying this to the real and imaginary parts shows
\[
 \Pr\left(\left|\sum_j\varepsilon_j\alpha_j\right|\geq u\right)
 \leq4e^{-qu^2/4}.
\]
There are $qs$ entries.  With $u=4\sqrt{(n+1)/q}$, the union-bound failure
probability is at most
\[
 4qs e^{-4(n+1)}\leq4\cdot2^{2n}e^{-4(n+1)}<1.
\]
We may therefore fix $E_2$ such that every entry of $C'$ has modulus at most
$4\sqrt{(n+1)/q}$.  Also $\norm{C'}\leq1$.

Let $h_i,c_i$ be row $i$ of $W,C'$, respectively, regarded as row vectors.
Then
\[
 \norm{h_i}^2=s/q,\qquad
 \norm{c_i}^2\leq v:=16s(n+1)/q.
\]
Choosing independent signs $\varepsilon_i$ for $E_1$, write
\[
 D=\sum_i\varepsilon_i h_i^*c_i,
 \qquad
 L_i=\begin{pmatrix}0&h_i^*c_i\\c_i^*h_i&0\end{pmatrix}.
\]
The two diagonal blocks of $\sum_iL_i^2$ satisfy
\[
 \sum_i\norm{c_i}^2h_i^*h_i\preceq vW^*W=vI_s,
 \qquad
 \sum_i\norm{h_i}^2c_i^*c_i
 =\frac sq C'^*C'\preceq\frac sq I_s\preceq vI_s.
\]
The norm of $\sum_i\varepsilon_iL_i$ equals $\norm D$.  Apply
Lemma~\ref{contract:moment} with $h=2s$ and $a=n+1$.  Since
$(2s)^{1/(2n+2)}\leq\sqrt2$, we obtain
\[
 \E\norm D
 \leq 2\sqrt2\sqrt{2(n+1)\,16s(n+1)/q}
 =16(n+1)\sqrt{s/q}
 \leq\frac14.
\]
The last inequality uses $s\leq q/M$.  Some choice of $E_1$ consequently
satisfies $\norm D\leq1/4$, which is more than required.

It remains to check the circuit structure.  A permutation on the $n$-bit
index register moves the last $q$ valid indices to the first $q$ valid
indices, an aligned $q$-element block.  On that block, $H_q$ uses at most
$n$ Hadamard slots, controlled on
the remaining bits; each diagonal sign matrix uses one diagonal slot.  Undo
the permutation afterward.  Lemma~\ref{slot:permutation} supplies clean
$O(n)$-slot implementations of the permutations.  All operations preserve
the valid $d$-dimensional subspace and are extended unitarily to unused index
states.  Labels change only the row permutation or diagonal-sign tables,
so the same slot structure works for every $U_\lambda$.
\end{proof}

The random signs are used only to prove that suitable tables exist.
Each label may use its own choices; the resulting slot circuit is deterministic.
The size of the removed block is large enough to make recursion short.
\begin{lemma}[Dimension decrease]\label{contract:depth}
For the parameters in~\eqref{contract:orders}, $s\geq d/(4M)$.
Starting at any $d\leq2^n$ and repeatedly replacing $d\geq2$ by $d-s$, one
reaches $d=1$ in fewer than
\begin{equation}\label{contract:J}
 J=8M(n+1)
\end{equation}
steps.
\end{lemma}

\begin{proof}
If $q/M\geq2$, then
$s=\lfloor q/M\rfloor\geq q/(2M)>d/(4M)$, because $d<2q$.
If $q/M<2$, then $s=1$ and $d<2q<4M$, with the same conclusion.
Thus every nontrivial step multiplies the order by at most $1-1/(4M)$.
If there were $J$ such steps, their final order would be at most
\[
 2^n(1-1/(4M))^J
 \leq2^n e^{-J/(4M)}
 =2^n e^{-2(n+1)}<1,
\]
contrary to its being a positive integer.
\end{proof}
